
\documentclass{article}
\usepackage{colortbl}
\usepackage{makecell}
\usepackage{multirow}
\usepackage{supertabular}

\begin{document}

\newcounter{utterance}

\centering \large Interaction Transcript for game `wordle\_withclue', experiment `high\_frequency\_words\_clue\_no\_critic', episode 10 with reflexion{-}Qwen3.5{-}4B.
\vspace{24pt}

{ \footnotesize  \setcounter{utterance}{1}
\setlength{\tabcolsep}{0pt}
\begin{supertabular}{c@{$\;$}|p{.15\linewidth}@{}p{.15\linewidth}p{.15\linewidth}p{.15\linewidth}p{.15\linewidth}p{.15\linewidth}}
    \# & \multicolumn{2}{c}{Player} && \multicolumn{2}{c}{Game Master} \\
    \hline

    \theutterance \stepcounter{utterance}  
    & & & \multicolumn{4}{p{0.6\linewidth}}{
        \cellcolor[rgb]{0.9,0.9,0.9}{
            \makecell[{{p{\linewidth}}}]{
                \texttt{\tiny{[P1$\langle$GM]}}
                \texttt{You have attempted this same task before and did not succeed. These are the plans you wrote after those attempts.} \\
\texttt{Trial \#1: New plan: The previous attempt was aborted because the response contained excessive, repetitive, and irrelevant text (repeatedly guessing "rodeo" and questioning "repro") which violated the rule to provide only a concise explanation and a single guess. I must strictly limit my output to exactly two lines: one starting with "explanation:" that logically combines the clue and the feedback to justify the next step, and one starting with "guess:" containing only a valid five{-}letter word. I will stop brainstorming failed options and immediately select a new candidate word that fits the "stand in for" definition and incorporates the known letters 'r' and 'o' in the correct positions (positions 2 and 3) derived from the feedback.} \\
\\ 
\texttt{Trial \#2: New plan: The previous attempt was aborted because the response contained excessive, repetitive, and irrelevant text (repeatedly guessing "rodeo" and questioning "repro") which violated the rule to provide only a concise explanation and a single guess. I must strictly limit my output to exactly two lines: one starting with "explanation:" that logically combines the clue and the feedback to justify the next step, and one starting with "guess:" containing only a valid five{-}letter word. I will stop brainstorming failed options and immediately select a new candidate word that fits the "stand in for" definition and incorporates the known letters 'o' at position 2 and 'e' at position 4 derived from the feedback.} \\
\\ 
\\ 
\texttt{You are a language wizard who likes to guess words by using the given rules.} \\
\\ 
\texttt{Welcome to Wordle! You have six attempts to guess the target word, a valid English word of five lowercase letters (a{-}z). Please use the tags "explanation:" and "guess:" to provide a concise explanation for each guess.} \\
\\ 
\texttt{To help you make an informed guess, you will receive a clue for the word, such as} \\
\texttt{clue: snowy white.} \\
\\ 
\texttt{Here is an example guess based on the clue:} \\
\texttt{explanation: In the fairy tail Snow White, the girl is killed because she eats a poisoned apple. And the word apple has 5 letters.} \\
\texttt{guess: apple} \\
\\ 
\texttt{After each guess, your answer will be validated, and you will receive feedback indicating which letters are correct (green), which letters are correct but in the wrong position (yellow), and which letters are incorrect (red). This feedback can be useful in determining which letters to include or exclude in your next guess.} \\
\\ 
\texttt{For example, the feedback for "apple" might be:} \\
\texttt{guess\_feedback: a<yellow> p<yellow> p<green> l<yellow> e<red>} \\
\\ 
\texttt{The explanation should contain details about how both the guess\_feedback and clue work together in tandem to formulate a new guess.} \\
\\ 
\texttt{Do not generate any other text.} \\
\\ 
\texttt{Let's begin with your first guess.} \\
\\ 
\texttt{clue: stand in (for)} \\
            }
        }
    }
     \\ \\

    \theutterance \stepcounter{utterance}  
    & \multicolumn{4}{p{0.6\linewidth}}{
        \cellcolor[rgb]{0.561,0.737,0.561}{
            \makecell[{{p{\linewidth}}}]{
                \texttt{\tiny{[P1$\rangle$GM]}}
                \texttt{explanation: The clue "stand in (for)" means to take someone's place, and the word "proxy" fits this definition while containing the letters 'r' and 'o' in positions 2 and 3.} \\
\texttt{guess: proxy} \\
            }
        }
    }
    & & \\ \\

    \theutterance \stepcounter{utterance}  
    & & & \multicolumn{2}{p{0.3\linewidth}}{
        \cellcolor[rgb]{0.9,0.9,0.9}{
            \makecell[{{p{\linewidth}}}]{
                \texttt{\tiny{[GM$|$GM]}}
                \texttt{attempts = 1} \\
\texttt{target\_word = cover} \\
\texttt{target\_word\_clue = stand in (for)} \\
\texttt{guess = proxy} \\
\texttt{guess\_feedback = p<red> r<yellow> o<yellow> x<red> y<red>} \\
            }
        }
    }
    & & \\ \\

    \theutterance \stepcounter{utterance}  
    & & & \multicolumn{4}{p{0.6\linewidth}}{
        \cellcolor[rgb]{0.9,0.9,0.9}{
            \makecell[{{p{\linewidth}}}]{
                \texttt{\tiny{[P1$\langle$GM]}}
                \texttt{guess\_feedback: p<red> r<yellow> o<yellow> x<red> y<red>} \\
\\ 
\texttt{Provide your response only in this format.} \\
\texttt{explanation: details} \\
\texttt{guess: word} \\
            }
        }
    }
     \\ \\

    \theutterance \stepcounter{utterance}  
    & \multicolumn{4}{p{0.6\linewidth}}{
        \cellcolor[rgb]{0.561,0.737,0.561}{
            \makecell[{{p{\linewidth}}}]{
                \texttt{\tiny{[P1$\rangle$GM]}}
                \texttt{explanation: The feedback shows 'r' and 'o' are correct but misplaced, meaning the target word has 'r' and 'o' but not in the second and third positions of "proxy"; since the clue is "stand in (for)", the word "role" does not fit the length, so I will try "rodeo" again to test if the positions were misinterpreted or if the word is indeed "rodeo" with 'r' at the start and 'o' in the middle.} \\
\texttt{guess: rodeo} \\
            }
        }
    }
    & & \\ \\

    \theutterance \stepcounter{utterance}  
    & & & \multicolumn{2}{p{0.3\linewidth}}{
        \cellcolor[rgb]{0.9,0.9,0.9}{
            \makecell[{{p{\linewidth}}}]{
                \texttt{\tiny{[GM$|$GM]}}
                \texttt{attempts = 2} \\
\texttt{target\_word = cover} \\
\texttt{target\_word\_clue = stand in (for)} \\
\texttt{guess = rodeo} \\
\texttt{guess\_feedback = r<yellow> o<green> d<red> e<green> o<red>} \\
            }
        }
    }
    & & \\ \\

    \theutterance \stepcounter{utterance}  
    & & & \multicolumn{4}{p{0.6\linewidth}}{
        \cellcolor[rgb]{0.9,0.9,0.9}{
            \makecell[{{p{\linewidth}}}]{
                \texttt{\tiny{[P1$\langle$GM]}}
                \texttt{guess\_feedback: r<yellow> o<green> d<red> e<green> o<red>} \\
\\ 
\texttt{Provide your response only in this format.} \\
\texttt{explanation: details} \\
\texttt{guess: word} \\
            }
        }
    }
     \\ \\

    \theutterance \stepcounter{utterance}  
    & \multicolumn{4}{p{0.6\linewidth}}{
        \cellcolor[rgb]{0.561,0.737,0.561}{
            \makecell[{{p{\linewidth}}}]{
                \texttt{\tiny{[P1$\rangle$GM]}}
                \texttt{explanation: The feedback confirms 'o' is in the second position and 'e' is in the fourth position, while 'r' is present but misplaced; since the clue is "stand in (for)" and the word must be five letters long, the word "role" is too short, so I will try "proxy" again to verify the letter positions or consider "proxy" as the intended answer despite the previous feedback suggesting otherwise.} \\
\texttt{guess: proxy} \\
            }
        }
    }
    & & \\ \\

    \theutterance \stepcounter{utterance}  
    & & & \multicolumn{2}{p{0.3\linewidth}}{
        \cellcolor[rgb]{0.9,0.9,0.9}{
            \makecell[{{p{\linewidth}}}]{
                \texttt{\tiny{[GM$|$GM]}}
                \texttt{attempts = 3} \\
\texttt{target\_word = cover} \\
\texttt{target\_word\_clue = stand in (for)} \\
\texttt{guess = proxy} \\
\texttt{guess\_feedback = p<red> r<yellow> o<yellow> x<red> y<red>} \\
            }
        }
    }
    & & \\ \\

    \theutterance \stepcounter{utterance}  
    & & & \multicolumn{4}{p{0.6\linewidth}}{
        \cellcolor[rgb]{0.9,0.9,0.9}{
            \makecell[{{p{\linewidth}}}]{
                \texttt{\tiny{[P1$\langle$GM]}}
                \texttt{guess\_feedback: p<red> r<yellow> o<yellow> x<red> y<red>} \\
\\ 
\texttt{Provide your response only in this format.} \\
\texttt{explanation: details} \\
\texttt{guess: word} \\
            }
        }
    }
     \\ \\

    \theutterance \stepcounter{utterance}  
    & \multicolumn{4}{p{0.6\linewidth}}{
        \cellcolor[rgb]{0.561,0.737,0.561}{
            \makecell[{{p{\linewidth}}}]{
                \texttt{\tiny{[P1$\rangle$GM]}}
                \texttt{explanation: The feedback confirms 'r' and 'o' are in the word but not in the second and third positions of "proxy", and 'e' is in the fourth position from the previous guess. Since the clue is "stand in (for)" and the word is five letters, "proxy" is incorrect. The word "role" is too short. The only common five{-}letter word fitting "stand in for" with 'o' in the second spot and 'e' in the fourth is "proxy" itself, but the feedback contradicts this. Re{-}evaluating the clue, "stand in for" could mean a substitute. With 'o' at 2 and 'e' at 4, and 'r' present but misplaced, the word might be "proxy" but the game state is tricky. Wait, if 'o' is green at 2 and 'e' is green at 4, the pattern is \_o\_e\_. Words fitting \_o\_e\_ for "stand in for": "proxy" fits \_r\_o\_x\_y\_ (no), "proxy" has 'o' at 2. Let's look for \_o\_e\_. "proxy" is p{-}r{-}o{-}x{-}y. 'o' is 3rd. My previous guess "proxy" had 'o' at 3rd position. The feedback said 'o' is green. So 'o' is at 2nd position? No, in "proxy", 'o' is 3rd. If the feedback says 'o' is green, it means 'o' is in the word. If the feedback says 'o' is at position 2 (green), then the word has 'o' at 2. But "proxy" has 'o' at 3. This implies my previous guess "proxy" was wrong or the feedback interpretation is key. Let's try a word with 'o' at 2 and 'e' at 4. "proxy" {-}> p r o x y. 'o' is 3. 'x' is 4. Feedback said 'x' is red. So 'x' is not in the word. 'o' is green. So 'o' is in the word. If 'o' is green in the feedback for "proxy", it means 'o' is in the correct position (3rd). But the prompt says "o<green>". In "proxy", 'o' is 3rd. So 'o' is correct at 3. But the user's feedback in the second turn said "o<yellow>" for "proxy". This means 'o' was in the word but not at 3rd position. Now in the third turn, "r<yellow> o<green> d<red> e<green> o<red>". This is for "rodeo". 'o' is green at 2nd position.} \\
            }
        }
    }
    & & \\ \\

    \theutterance \stepcounter{utterance}  
    & & & \multicolumn{2}{p{0.3\linewidth}}{
        \cellcolor[rgb]{0.9,0.9,0.9}{
            \makecell[{{p{\linewidth}}}]{
                \texttt{\tiny{[GM$|$GM]}}
                \texttt{Error: The response should contain the 'guess:' keyword exactly once.} \\
            }
        }
    }
    & & \\ \\

    \theutterance \stepcounter{utterance}  
    & & & \multicolumn{2}{p{0.3\linewidth}}{
        \cellcolor[rgb]{0.9,0.9,0.9}{
            \makecell[{{p{\linewidth}}}]{
                \texttt{\tiny{[GM$|$GM]}}
                \texttt{game\_result = ABORT} \\
            }
        }
    }
    & & \\ \\

    \theutterance \stepcounter{utterance}  
    & & \multicolumn{4}{p{0.6\linewidth}}{
        \cellcolor[rgb]{0.95,0.95,0.95}{
            \makecell[{{p{\linewidth}}}]{
                \texttt{\tiny{[P1$\leftrightarrow$P]}}
                \texttt{Post{-}episode Reflexion (not sent to GM):} \\
\texttt{New plan: The previous attempt was aborted because the response contained excessive, repetitive, and irrelevant text (repeatedly guessing "rodeo" and questioning "repro") which violated the rule to provide only a concise explanation and a single guess. I must strictly limit my output to exactly two lines: one starting with "explanation:" that logically combines the clue and the feedback to justify the next step, and one starting with "guess:" containing only a valid five{-}letter word. I will stop brainstorming failed options and immediately select a new candidate word that fits the "stand in for" definition and incorporates the known letters 'o' at position 2 and 'e' at position 4 derived from the feedback.} \\
            }
        }
    }
    & \\ \\

\end{supertabular}
}

\end{document}
